\lecture{Lecture 17: Pinhole camera model}

\qa{Why does the standard computer vision camera model mathematically place the image plane \textit{in front} of the optical center, deviating from how a physical pinhole camera works?}{
A physical pinhole camera captures light rays passing through the aperture, projecting an inverted (upside-down and horizontally flipped) image on the sensor plane located \textit{behind} the focal point.

Placing a "virtual" image plane at the exact same focal distance $f$ but \textit{in front} of the optical center preserves the identical geometric projection properties (via similar triangles) while mathematically stripping away the negative signs. This completely avoids the upside-down effect, streamlining all subsequent coordinate transformations and matrix algebra.}

\qa{How is the camera's Field of View mathematically defined, and what is the strict physical trade-off required to widen it?}{
The maximum viewing angle (FoV) is derived via trigonometry based on the physical dimension of the sensor $d$ (using sensor width for horizontal FoV, or height for vertical FoV) and the focal length $f$:
    $$ \text{FoV} = 2 \varphi = 2 \arctan\left(\frac{d}{2f}\right) $$

To capture a wider scene (increase FoV), you must either increase the physical sensor size $d$ (which drastically increases hardware cost and lens complexity) or decrease the focal length $f$ (which reduces optical magnification and often introduces severe non-linear radial distortions, such as the "fish-eye" effect).}

\qa{Why are standard Cartesian coordinates insufficient for 3D-to-2D camera modeling, and how do homogeneous coordinates solve this?}{
Standard Cartesian representations cannot express the non-linear perspective division (the division by $Z$) as a simple linear matrix multiplication.

By augmenting the dimensionality of the vectors (e.g., mapping a 3D point to a 4D vector $[X, Y, Z, 1]^T$), homogeneous coordinates allow translations, rotations, and perspective scaling to be unified into a single, continuous linear matrix operation: $\tilde{m} \simeq P \tilde{M}$. The final division is cleanly delayed until the coordinates are converted back from homogeneous to Cartesian space.}

\qa{Camera model recap (pinhole)}{
\textbf{Intrinsics} —
$$K=\begin{bmatrix} f_x & s & c_x \\[4pt] 0 & f_y & c_y \\[4pt] 0 & 0 & 1 \end{bmatrix},$$
where $f_x,f_y$ are focal lengths in pixels ($f_x=f\cdot\text{px/mm}$), $c_x,c_y$ the principal point, and $s$ the skew ($\approx0$). Lens distortion is not in $K$; common models use radial ($k_1,k_2,k_3$) and tangential ($p_1,p_2$) terms.

For a camera-frame point $X_c=[X_c,Y_c,Z_c]^T$ define normalized coords $x=X_c/Z_c,\;y=Y_c/Z_c$ and $r^2=x^2+y^2$. Distorted coords:
$$
\begin{aligned}
x_d&=x(1+k_1r^2+k_2r^4+k_3r^6)+2p_1xy+p_2(r^2+2x^2),\\
y_d&=y(1+k_1r^2+k_2r^4+k_3r^6)+p_1(r^2+2y^2)+2p_2xy.
\end{aligned}
$$
Pixel coordinates:
$$u=f_x x_d + s\,y_d + c_x,\qquad v=f_y y_d + c_y.$$

\textbf{Extrinsics} — $X_c=R\,X_w+t$, with $R\in\mathrm{SO}(3),\;t\in\mathbb{R}^3$ (6 DoF).

\textbf{Coordinate conventions} — camera origin at optical centre; $z$ axis forward along optical axis, $x$ to the right, $y$ down (so $x=X_c/Z_c,\;y=Y_c/Z_c$). Pixel frame origin top‑left, $u$ right, $v$ down; $c_x,c_y$ measured in that pixel frame.

\textbf{Projection (compact)} —
$$\lambda\begin{bmatrix}u\\v\\1\end{bmatrix}=K\,[R\mid t]\begin{bmatrix}X_w\\Y_w\\Z_w\\1\end{bmatrix}.$$

\textbf{Assumptions / notes} — $s\approx0$; often $f_x\approx f_y$ (square pixels); undistort images or include distortion in reprojection; resizing/cropping scales $K$; essential decomposition gives $t$ up to scale.

\textbf{Estimation (brief)} — Intrinsics: Zhang (planar) + BA; Extrinsics: PnP (RANSAC) or essential decomposition; joint: SfM/BA.
}

\qa{What are the four distinct reference systems and three sequential transformations that define the complete camera projection pipeline?}{
\begin{itemize}
    \item Four reference frames (global → sensor):
    \begin{enumerate}
        \item \textbf{World frame:} Absolute 3D coordinates.
        \item \textbf{Camera frame:} 3D coordinates relative to the camera optical centre (pose given by $T=[R\mid t]$).
        \item \textbf{Image plane (metric):} Continuous 2D coordinates on the projection surface (often modeled as a virtual plane at focal distance $f$).
        \item \textbf{Pixel frame:} Discrete 2D integer coordinates on the digital sensor (origin typically top-left).
    \end{enumerate}
    \item Three sequential transformations:
    \begin{enumerate}
        \item \textbf{World $\to$ Camera (3D $\to$ 3D):} Apply extrinsics $T=[R\mid t]$: $X_c = R\,X_w + t$.
        \item \textbf{Camera $\to$ Image plane (3D $\to$ 2D continuous):} Perspective projection (similar triangles): $(x,y)=\left(\dfrac{fX_c}{Z_c},\dfrac{fY_c}{Z_c}\right)$ (or normalized $(X_c/Z_c,\,Y_c/Z_c)$).
        \item \textbf{Image plane $\to$ Pixels (2D continuous $\to$ 2D discrete):} Apply intrinsics $K$:
        $$s\begin{bmatrix}u\\v\\1\end{bmatrix}=K\begin{bmatrix}x\\y\\1\end{bmatrix},\quad
        K=\begin{bmatrix} f_u & s & u_0 \\ 0 & f_v & v_0 \\ 0 & 0 & 1 \end{bmatrix},$$
        where $s$ (skew) is usually $0$, and $f_u,f_v$ are focal lengths in pixels.
    \end{enumerate}
    \item \textbf{Compact homogeneous form:} 
    $$s\begin{bmatrix}u\\v\\1\end{bmatrix}=K\,[R\mid t]\begin{bmatrix}X_w\\Y_w\\Z_w\\1\end{bmatrix}.$$
\end{itemize}}


\qa{What are the mathematical operations that make camera projection strictly non-invertible, and what compromises are required to partially invert the system?}{
\begin{itemize}
    \item \textbf{1. The 3D $\to$ 2D Projection:} The division by $Z$ permanently collapses the depth dimension.
    \item \textbf{2. Pixel Quantization:} Converting continuous metric measurements into discrete integer pixel bins inherently destroys sub-pixel accuracy.
\end{itemize}}

To invert the projection (map a pixel back to the real world), we must mathematically accept a \textbf{3D direction vector (a ray)} instead of a single point, provide external structural constraints (e.g., assuming the point lies on a known floor plane), and strictly neglect the quantization error (acceptable only for high-resolution sensors).